\chapter{Conclusion}
\label{chap:conclusion}

\noindent
Automated intrusion response is still rarely enabled on operational technology, and the reason is trust rather than detection: an operator will not arm a defence that might cut the process it is meant to protect (Chapter~\ref{chap:introduction}). This project set out to lower that barrier with CARS, an SDN intrusion-response system whose action is conditioned on the criticality of the asset and the industrial operation being attempted, and whose responses are bounded in time, reversible, and evidence-generating. It was built on a hardware testbed with two real Siemens S7-1212C PLCs and a live tank process, and evaluated at the level of individual packets. The central result is that the armed defence refused every illegitimate operation while cutting none of the legitimate process traffic: a false-positive rate of zero on the live process, with a response that scales by criticality and enforcement shown on the wire rather than inferred from a log. That property, more than raw detection, is what an operator needs before arming a reactive defence on a running plant.

\section{Contributions}
\label{sec:contributions}

The project met the five objectives set out in Chapter~\ref{chap:introduction}, summarised against their outcomes in Table~\ref{tab:aims}. A representative three-node SDN-ICS testbed was built with real PLCs, an operator panel, an engineering workstation and a hardware-in-the-loop tank, so that every response could be judged against a physical consequence rather than a simulated one. On that fabric, CARS was realised as a three-table OpenFlow pipeline, identity binding against spoofing, a stateful policy stage carrying the criticality- and operation-aware rulebook, and an L2 switch, behind which sit a graded seven-level response ladder, a flow-integrity self-check and an authenticated control interface. The decision logic was measured across the range of source role, operation class and asset criticality and returned a correct verdict for every case, with no legitimate operation wrongly cut, and the same verdicts were corroborated by the live campaign. Each capability was then validated at the packet level, protected against unprotected, from reconnaissance through identity spoofing, out-of-state traffic, control-plane and flow tampering, operation-aware commands, and a false-data injection that overflowed the tank unprotected yet was severed before a single write reached the PLC when the defence was armed. The reaction window was characterised directly, with a median mitigation of 12.6\,ms, and the whole account was drawn together in an end-to-end trace of two attacker positions from one machine, which showed the same enforcement applied at coarse gateway granularity for a NAT-collapsed external attacker and at fine per-host granularity for an on-segment insider.

\begin{table}[t]
\centering
\small
\begin{tabular}{p{0.46\textwidth}p{0.46\textwidth}}
\hline
Objective & Outcome \\
\hline
Build a representative hardware testbed with an SDN fabric, real PLCs, an operator panel and a live process & Three-node fabric, two S7-1212C PLCs, KTP700 panel, EWS and a Factory IO tank driven through the controlled switches (Chapter~\ref{chap:execution}) \\
Design a graded, bounded, reversible, evidence-generating response model & Seven-rung ladder (ALLOW to REFUSE) with criticality-scaled self-healing durations (75/60/45/30\,s) and an audit record per decision (Chapters~\ref{chap:execution},~\ref{chap:evaluation}) \\
Implement it as an OpenFlow controller with a layered pipeline and integrity checks & Three-table pipeline (GUARD, stateful POLICY, SWITCH), flow-integrity checker and token-authenticated control API (Chapters~\ref{chap:execution},~\ref{chap:evaluation}) \\
Integrate a hardware-in-the-loop process for real-consequence testing & Factory IO tank on the live PLC, judged by overflow versus safe band (Chapters~\ref{chap:execution},~\ref{chap:evaluation}) \\
Evaluate accuracy and false positives; validate each capability at the packet level, protected versus unprotected & 100\% decision accuracy, 0\% false positives; wire-level campaign across every capability; reaction window characterised (median 12.6\,ms) (Chapter~\ref{chap:evaluation}) \\
\hline
\end{tabular}
\caption{Aims and objectives against outcomes.}
\label{tab:aims}
\end{table}

\section{What the evaluation establishes, and what it does not}
\label{sec:standing}

The result is best read as a layered account of responsibility rather than a single defence. The proactive layer, identity binding and a stateful default-deny allowlist, stops a source that should never connect with no reaction window at all, which the two-pivot trace confirmed for both the external and the on-segment attacker. The reactive, operation-aware layer governs what an already-permitted party may then do, refusing a dangerous operation on a critical asset while permitting a read; the volumetric overlay cuts high-rate abuse even from a trusted role; and, added as defence in depth, a read-only process guardian raises an alarm on the physical anomalies the network layer must not act on. Together these draw a clean division: the network decides who and what may connect, and a process-invariant monitor watches the consequence.

The evaluation is honest about where a lab prototype stops. It is a functional-prototype validation on a single testbed with one hardware-in-the-loop process, so the zero-false-positive figure is a statement about this system and its rulebook across the decision space and the live campaign, not a statistical guarantee for arbitrary industrial traffic. The reaction window is a few milliseconds here but is a property of any detect-then-respond defence, and a faster process would demand more of the proactive layer. The residual adversary the design cannot fully answer is a trusted party manipulating the process slowly and within its physical envelope; the guardian narrows this case but does not close it, because prevention on the safety loop is the very thing that must not be automated. Several further limits are architectural rather than flaws in the response model, and are carried openly in Chapter~\ref{chap:evaluation}: the OpenFlow control channel runs in plaintext on an isolated management plane, the flow-integrity checker polls on an interval that a transient injection can beat, the fabric was not tested against a spoofed-source state-exhaustion flood, and the file-based detection bridge would bottleneck under a heavy distributed load.

\section{Future work}
\label{sec:futurework}

The limitations above map directly onto the next steps.

\begin{itemize}
\item \textbf{Toward deployment.} Run OpenFlow over TLS and measure the added control-path latency; replace the ten-second flow-integrity poll with an event-driven flow-monitor subscription so a transient rule change cannot slip between polls; test and harden the fabric against spoofed-source state-exhaustion flooding with connection-tracking limits and per-source metering; and move the detection bridge from file-and-HTTP to in-memory inter-process communication to hold the reaction window under heavy load.
\item \textbf{Deepening the process layer.} Generalise the process guardian into sensor-aware remediation on the live process, with setpoint- and invariant-level anomaly detection, so that a trusted insider driving the process slowly within its envelope can be flagged rather than only a gross excursion.
\item \textbf{Extending the response repertoire.} Give the DEFLECT path a fully interactive decoy that sustains a session with the attacker, and add a TCP reset on quarantine so a recoverable endpoint fails fast instead of hanging on a silent drop.
\item \textbf{Broader validation.} Exercise the decision logic against a larger and noisier traffic sample, more assets and alternative rulebooks, and a faster process, to move the accuracy result from a functional-prototype figure toward a statistical one; and extend enforcement across the full multi-cell fabric.
\item \textbf{Advanced deception.} Explore replica-based moving-target defence as an alternative to block-and-maintain, trading a decoy and a shifting target for a graded cut \cite{etxezarreta2026replica}.
\end{itemize}

\noindent
CARS does not claim to solve intrusion response for industrial control systems. What it shows is narrower and, for the deployment question, more useful: that a reactive SDN defence can be made criticality- and operation-aware, bounded and reversible enough to stop real attacks on a live process at the packet level without ever cutting the process itself. On the evidence of this testbed, the trust barrier that keeps such defences switched off is not fundamental, and an operator given a defence that refuses an unsafe action before it will interrupt the plant has reason to arm it.
